% Converted from sections/02-related-work.md by tools/md2tex.py; edit this file from now on.
\section{Related work}\label{sec:2}

Four lines of work bear on this report. They are reviewed from the most general to the closest:
declarative configuration and its convergence, reconciliation in cluster managers, bridging
declarative models and imperative interfaces, and the analysis and generation of tooling from API
descriptions. The section ends with a summary of how the present method differs.

\subsection{Declarative configuration and convergence}\label{sec:2.1}

Declarative configuration management predates cloud IaC. CFEngine introduced \emph{convergent} operations,
which move a host towards the state a policy implies whatever its starting point
\citep{burgess2005}. Couch and Sun formalised convergence
algebraically \citep{couch2003}. They also showed
that, for a single host and some classes of process, reproducibility of the observed effect can be
verified statically \citep{couch2004}. Testing work
turned these notions into checks on real scripts. \citet{hummer2013}
proposed model-based testing of idempotence and found non-idempotent automations in public
repositories. \citet{hanappi2016} tested \emph{reliable
convergence} under transient faults. They showed that when a pairwise \emph{preservation} property holds
between resources, the whole configuration converges without every execution order being enumerated.
Rehearsal checks Puppet manifests for determinism and idempotence
\citep{shambaugh2016}. Tortoise repairs configurations
after manual changes have made a system drift from its specification
\citep{weiss2017}. More generally, idempotence is the
property that makes retry-based fault tolerance correct
\citep{ramalingam2013}.

Empirical studies of IaC have since characterised what goes wrong in practice. These include defect
taxonomies of IaC scripts \citep{rahman2020}, security
smells \citep{rahman2019b}, and faults in the engines of
Ansible, Puppet and Chef \citep{drosos2024}.
\citet{hassan2024} studied 5,110 defects in \emph{state
reconciliation}, the step that queries and compares infrastructure state before changing it, and
found eight categories, with infrastructure inventory the most frequent. Tool research has addressed
program testing for general-purpose IaC languages
\citep{sokolowski2024} and semantic checks that
Terraform programs violate only at deployment time
\citep{qiu2024}. Other work has addressed bugs that
appear only when an existing deployment is updated
\citep{xiang2025} and lifting existing cloud state
into IaC programs \citep{peng2025}. Formal work has
begun to verify plans: BPlan models Terraform in Maude and checks that applying a plan never
duplicates a resource or leaves a dependency dangling
(\citealp{perdereau2026}; research artifact of a 2026
conference paper).

This literature supplies the vocabulary used here (convergence, idempotence, drift, reconciliation
defects) and several of the properties the evaluation tests. It studies engines and programs above
an existing provider layer, however, and does not address how that layer is built when the API offers
no resources.

\subsection{Reconciliation in cluster managers}\label{sec:2.2}

Borg, Omega and Kubernetes organise cluster management as a set of controllers that each observe
shared state and act to converge it on a declared target
\citep{burns2016}. This design has prompted a line of work on
controller correctness. Sieve perturbs a controller's view of cluster state (stale, intermediate and
unobserved states) and compares the resulting state evolution with an unperturbed run; it found 46
safety and liveness bugs in ten controllers
\citep{sun2022}. Acto tests the
operation correctness of Kubernetes operators end to end
\citep{gu2023}. Anvil verifies controller liveness
formally (Sun et al.~2024, OSDI, pp.~649--666). Empirical studies have catalogued operator bugs
\citep{xu2024} and conflicts between controllers that
reconcile the same objects \citep{an2026}.
Kubernetes also defines field ownership for declarative updates (server-side apply), which records
which manager owns which field \draftnote{cite the Kubernetes server-side apply documentation}. The present method adopts a related rule: a document owns exactly
the fields it names.

These systems reconcile against an API server that already stores declarative objects with a
separation between specification and status. Sieve's design depends explicitly on that
state-centric interface. A procedure-oriented API offers no such interface, so the separation
between desired and observed state has to be constructed on the client side.

\subsection{Bridging declarative models and imperative interfaces}\label{sec:2.3}

The closest earlier work connects declarative models with imperative or procedural interfaces.

\citet{maghraoui2006} bridged declarative object models
and procedural provisioning tools: AI planning generated the workflow that takes a data centre from
its current to its desired state. The CHAMPS system derived parallel change plans from dependency
information \citep{keller2004}. The TOSCA standard
describes applications declaratively as topologies and imperatively as management plans
\citep{binz2013}, and imperative plans can be
generated from declarative topologies
\citep{breitenbucher2014}.
\citet{wettinger2014} contrasted convergent
deployment automation with a compensation-based alternative built on snapshots. Theoretical results
bound what such planners can achieve. Plan achievability is undecidable in the general Aeolus
component model and tractable only under restrictions
\citep{cosmo2014}. Quiescence defines when a running
system can be changed safely \citep{kramer1990}.

Research on models at run time keeps a model causally connected to a running system, so that the
model reflects the system and editing the model changes it
(\citealp{blair2009};
\citealp{morin2009}; survey in
\citealp{bencomo2019}). Song and colleagues came closest
to the present problem:

\begin{itemize}
\item
  they generated synchronisation engines from a meta-model plus a description of how to manipulate
  its elements through the system's management API
  \citep{song2010a};
\item
  they inferred such meta-models from client code that calls the API
  \citep{song2010b};
\item
  they handled conflicts between model and system changes
  \citep{song2011}.
\end{itemize}

Related model-driven work built bridges between APIs and models
\citep{izquierdo2012}, extracted models from Web API
documentation \citep{rodriguez2012} and managed
cloud resources through run-time models
\citep{zhang2013}. In the cloud setting, CloudMF
combines a provisioning language with a models-at-run-time engine that adapts a running multi-cloud
deployment to a target model \citep{ferry2018}. Tooling built on the Open Cloud Computing
Interface (OCCI) generates connectors and compares topologies to adapt them at run time
\citep{zalila2019,korte2018,erbel2018,challita2021}.

Formally, the mapping between a declarative document and remote state is a view-update problem.
Lenses express it as a pair of transformations with round-trip laws
(\citealp{foster2007}; cross-disciplinary overview in
\citealp{czarnecki2009}). Three-way merging, which
planning against desired, recorded and observed state resembles, has weaker guarantees than often
assumed \citep{khanna2007}.

This body of work established that declarative intent can drive imperative interfaces through
generated or modelled mappings. Most of it, however, targets interfaces that are already
resource-oriented (OCCI is CRUD over resources by design) or relies on connectors written per
resource type. None of it provides a mechanical check that a mapping covers every field the
interface exposes, and the evaluations are largely case demonstrations rather than tests of
convergence and recovery.

\subsection{Analysing and generating tooling from API descriptions}\label{sec:2.4}

REST-testing research has developed techniques to recover the structure of an API from its
description. RESTler infers producer--consumer dependencies between requests from an OpenAPI document,
so that a request consuming an identifier is ordered after the request producing it
\citep{atlidakis2019}. Morest maintains a property graph
of such dependencies and repairs it during execution
\citep{liu2022}. A comparison of ten tools identified
the construction of valid operation sequences as the central difficulty
\citep{kim2022}. OpenAPI cannot express dependencies
between parameters, for which a dedicated language has been proposed
\citep{martinlopez2022}. Studies of API design have
assessed cloud REST APIs against design practice
\citep{petrillo2016} and collected structural
metrics over more than a million OpenAPI descriptions
\citep{serbout2024}.

On the tooling side, HashiCorp's OpenAPI provider generator turns an OpenAPI 3 document into a
Terraform provider code specification. A developer must still write a configuration that maps
operations to resources, because a general HTTP API description does not align with the provider
model on its own; the output covers schema and data handling, not request logic
(\href{https://github.com/hashicorp/terraform-plugin-codegen-openapi}{terraform-plugin-codegen-openapi}).
Crossplane's Upjet (\href{https://github.com/crossplane/upjet}{github.com/crossplane/upjet}) and Pulumi's
Terraform bridge (\href{https://github.com/pulumi/pulumi-terraform-bridge}{github.com/pulumi/pulumi-terraform-bridge})
generate providers from existing Terraform providers, so they inherit resource semantics rather
than derive them. AWS Cloud Control API takes the opposite route: the vendor publishes, for each
resource type, a schema that declares the primary identifier and the read-only, create-only and
write-only properties, and a uniform create--read--update--delete--list interface follows from it \draftnote{cite the AWS Cloud Control API and CloudFormation resource-schema
documentation}.

These generators and schemas show which information a declarative tool needs from an API. When the
API is resource-oriented they recover much of it from the description; when it is not, they leave it
to hand-written configuration or code. The kind specification proposed here can be read as a
Cloud-Control-style resource schema written by the integrator for an API whose vendor did not
provide one. It is extended with what a procedure-oriented API additionally requires: the role of
each procedure, write groups, identity discovery, read authority, and structural shapes.

\subsection{Summary}\label{sec:2.5}

Table~\ref{tab:1} summarises the comparison. The present method combines elements that earlier work treated
separately:

\begin{itemize}
\item
  the convergence and recovery properties studied for configuration management and controllers;
\item
  the model-to-interface mapping studied for models at run time and cloud modelling;
\item
  the description-driven generation used in API tooling.
\end{itemize}

It adds what none of them provides for a procedure-oriented API: a per-kind specification of the
missing semantics, a mechanical check that the specifications cover every field the API exposes, and
conformance scenarios derived from the same specifications.

\begin{table}[tbp]
\centering
\small
\caption{Positioning of the method against related work.}
\label{tab:1}
\begin{tabularx}{\linewidth}{LLLLlL}
\toprule
Line of work & Representative work & Target interface & Mapping to the interface & Completeness check & Convergence and recovery evaluated \\
\midrule
Convergent configuration & \citealp{hanappi2016}; \citealp{shambaugh2016} & Host resources & Built into the tool's resource types & No & Yes \\
Cluster reconciliation & \citealp{sun2022}; \citealp{gu2023} & State-centric API server & Controller code per resource & No & Yes \\
Model-driven provisioning, TOSCA & \citealp{maghraoui2006}; \citealp{breitenbucher2014} & Scripts and services & Hand-modelled plans and artifacts & No & Limited \\
Models at run time & \citealp{song2010a,song2010b} & Management APIs & Generated from a meta-model and API description & No & Case studies \\
Cloud modelling, OCCI & \citealp{ferry2018}; \citealp{zalila2019} & Resource-oriented cloud APIs & Connectors per resource type & No & Case studies \\
API tooling generators & terraform-plugin-codegen-openapi; Upjet & OpenAPI or existing providers & Configuration or inherited code & No & No \\
\textbf{This report} & dokploy-iac & \textbf{Procedure-oriented RPC API} & \textbf{Declarative kind specifications, no per-kind hooks} & \textbf{Yes (coverage ledger)} & \textbf{\result{}} \\
\bottomrule
\end{tabularx}
\end{table}
